\chapter{OsdagBridge Overview and Setup}
\label{ch:setup}

Before diving into individual features, this chapter gives a map of the OsdagBridge codebase
and explains how a bridge design flows through the application. Understanding this structure
makes the rest of the report easier to follow, because each later chapter modifies a small,
well-defined part of this whole.

\section{Repository Structure}
OsdagBridge is a Python application organised as a single installable package under
\texttt{src/osdagbridge}. The repository root contains the usual packaging and environment
files (\texttt{pyproject.toml}, \texttt{requirements.txt}, \texttt{environment.yml}), plus
\texttt{docs/}, \texttt{tests/}, \texttt{examples/plate\_girder/}, and the source tree. The
important part for this report is the package itself:

\begin{lstlisting}[language=bash,caption={High-level layout of the OsdagBridge package.},label={lst:tree}]
src/osdagbridge/
    core/                     # backend: analysis, design, reports
        bridge_types/
            plate_girder/
                plategirderbridge.py    # plate-girder bridge module
                plot_generator.py       # 2D/3D plot generation
        bridge_components/    # reusable structural elements
        reports/
            report_generator.py         # design-report builder
            __init__.py
        input/  output/  solvers/  design_codes/
    desktop/                  # frontend: PySide6 GUI
        ui/
            cad_3d.py                   # 3D CAD scene assembly
            utils/
                custom_3dviewer.py      # OCC viewer + navigation modes
            docks/
                input_dock.py           # Input Dock
                output_dock.py          # Output Dock
                cad_top_view.py         # top / plan plot view
    cli/                      # command-line interface
    web/                      # Django + React web application
\end{lstlisting}

The two halves that matter most for my work are:
\begin{itemize}[leftmargin=1.4em]
  \item \textbf{\texttt{core/} --- the backend.} Pure-Python design logic with no GUI
        dependencies. The plate-girder bridge type
        (\texttt{core/bridge\_types/plate\_girder/plategirderbridge.py}) computes the design,
        and the \texttt{core/reports} package turns the results into a PDF design report.
  \item \textbf{\texttt{desktop/} --- the frontend.} A PySide6/Qt GUI. The
        \texttt{ui/docks} package holds the Input Dock, Output Dock, and plot views, while the
        CAD subsystem hosts the 3D viewer and its toolbar.
\end{itemize}

\section{The Plate-Girder Bridge Design Flow}
A user's interaction with OsdagBridge follows a consistent pipeline, and every feature in this
report plugs into one of its stages:

\begin{enumerate}[leftmargin=1.6em]
  \item \textbf{Input.} The user enters bridge parameters and material properties in the
        \emph{Input Dock}. Inputs are validated before a design can run
        (Chapter~\ref{ch:validators}).
  \item \textbf{Design.} On pressing \emph{Design}, the plate-girder module solves the bridge
        --- girders, deck slab, shear studs, bracings, and diaphragms --- and freezes the
        results into an output dictionary (Chapter~\ref{ch:geometry}).
  \item \textbf{Output \& visualisation.} Results are shown in the \emph{Output Dock}; the
        \emph{3D CAD window} renders the model with nodes, a grillage, and a coordinate axis,
        all navigable through a toolbar (Chapters~\ref{ch:cad}, \ref{ch:toolbar},
        \ref{ch:axis}); 2D plots are drawn in the plot views
        (Chapter~\ref{ch:plots}).
  \item \textbf{Reporting.} The user generates a detailed, IS-code-based \emph{design report}
        as a PDF, including calculations, summary tables, and plots
        (Chapters~\ref{ch:report}, \ref{ch:plots}).
\end{enumerate}

\osdagfig{figure2.1.png}{The OsdagBridge desktop application, showing the Input Dock,
the 3D CAD window with a plate-girder bridge, the Output Dock, and the message log.}{fig:gui}

\section{Development Setup}
OsdagBridge is developed as a conda-based Python environment. A typical contributor workflow,
and the one I used throughout the internship, is:

\begin{lstlisting}[language=bash,caption={Local development setup.},label={lst:setup}]
# clone and enter the repository
git clone https://github.com/Aditya-Donde/OsdagBridge.git
cd OsdagBridge

# create the environment
conda env create -f environment.yml
conda activate osdagbridge

# install the package in editable mode
pip install -e .

# launch the desktop application
python -m osdagbridge.desktop
\end{lstlisting}

Work was contributed through GitHub pull requests against the maintainer's branches
(\texttt{test}, \texttt{report}, and \texttt{dev}), each PR reviewed before being merged. The
chapters that follow are organised by feature area, in the order the work was carried out, and
each cites the pull request(s) in which the code was contributed.
